\documentclass[UTF8,fontset=fandol,11pt,a4paper]{ctexart}
\usepackage[margin=2.3cm]{geometry}
\usepackage{amsmath,amssymb,bm,booktabs,graphicx,array}
\usepackage[authoryear,round]{natbib}
\usepackage{caption,float,placeins}
\usepackage[colorlinks=true,citecolor=blue,linkcolor=blue,urlcolor=blue]{hyperref}
\usepackage{microtype}
\captionsetup{font=small,labelfont=bf}
\setlength{\parindent}{2em}
\setlength{\parskip}{0pt}
\linespread{1.12}
\newcommand{\fig}[3]{\begin{figure}[!htbp]\centering\includegraphics[width=#2\textwidth]{../Figures/#1.pdf}\caption{#3}\end{figure}}
\title{A2：多岭参数模型平均的下界独立复验与B任务预实验}
\author{第3组\quad 洪心驿}
\date{2026年10月22日}
\begin{document}
\maketitle
\begin{abstract}
第一次周报的A2实验发现，接近插值阈值时，包含极小岭参数的广义交叉验证模型平均（GCV-MA）可能低估样本外风险。本报告按第一次周报预定协议，以未参与原实验的新随机种子独立复验两个维数比、五个候选下界，共运行10个配置、每配置150次重复。两种维数比下，提高下界均降低GCV-MA超额风险，并使平均GCV分数与平均总风险之比趋近1。在此基础上仅提出一个B研究问题：训练数据内的折外评价能否稳定选择GCV-MA的候选下界，使其在插值附近兼顾风险校准与预测风险。另以独立种子运行200次预实验；五折选界明显优于完整网格，接近固定下界1，单次留出选界则不稳定。这些均为指定合成设置下的有限样本结果。
\end{abstract}

\section{研究衔接与问题设定}
\citet{yu}提出对不同岭参数估计作凸平均，并以GCV型准则选权；\citet{dw}研究高维比例增长下岭回归预测风险。第一次A2报告在相关设计及 $p/n\approx1$ 的探索实验中发现GCV-MA风险校准失真，第一次周报据此约定独立复验五个候选下界。这里仅完成洪心驿负责的A2任务与由其结果导出的B预实验；A1的新随机种子复验属于另一位组员的任务。

给定训练数据 $X\in\mathbb R^{n\times p},y\in\mathbb R^n$，第 $m$ 个候选为
\begin{equation}
 \widehat{\bm\beta}_{\lambda_m}=(X^\top X+n\lambda_m I)^{-1}X^\top y,
 \qquad \widehat{\bm\beta}(w)=\sum_{m=1}^{21}w_m\widehat{\bm\beta}_{\lambda_m},
 \quad w_m\ge0,\ \sum_mw_m=1.
\end{equation}
候选 $\lambda_m$ 在 $[\lambda_{\min},100]$ 上对数等距。GCV-MA最小化混合平滑器的GCV分数。主要评价量是条件超额预测风险
$L=(\widehat{\bm\beta}-\bm\beta)^\top\Sigma(\widehat{\bm\beta}-\bm\beta)$；总风险为 $L+1$。校准比定义为重复平均GCV分数除以重复平均总风险，接近1表示\emph{均值层面}校准，不能推出单次实验或整个混合平滑器族上的一致性。

\section{独立复验协议与可核验产物}
固定 $n=100$、Toeplitz协方差 $\Sigma_{jk}=0.7^{|j-k|}$、稠密高斯信号、信噪比1、同方差高斯噪声；令 $p/n\in\{1,1.5\}$，$\lambda_{\min}\in\{10^{-4},10^{-3},10^{-2},10^{-1},1\}$。每个配置150次重复；同一维数比的五个下界共享150份数据和折分，维数比之间的数据独立。新种子根为20261022，第一次报告使用20261007。保留原A2实现的样本量归一化、GCV凸优化、Oracle定义及评价协议。完整原始记录、候选权重、配置、执行环境和校验结果存于\texttt{Results/}，重跑入口为\texttt{Code/run\_a2.py}。

对照方法为同网格选参的Ridge-GCV、留一凸平均LOO-MA、以及使用真实 $\bm\beta,\Sigma$ 的Oracle-MA。最后一种仅刻画候选集合可达到的风险，不是可实施算法。报告表格中的标准误为150次独立数据重复的样本标准差除以 $\sqrt{150}$；五个下界之间比较使用相同重复编号的配对差值。模拟中真实参数只用于评价和Oracle参照，不参与可实施方法的调参。

\section{A2复验结果}
\begin{table}[!htbp]\centering\small
\caption{独立复验的平均超额风险与GCV-MA校准比。括号为GCV-MA均值标准误；其余三列为方法均值。}
\setlength{\tabcolsep}{4pt}
\begin{tabular}{ccccccc}\toprule
$p/n$ & 下界 & GCV-MA & Ridge-GCV & LOO-MA & Oracle-MA & 校准比\\\midrule
1 & $10^{-4}$ & 5.713 (0.564) & 0.895 & 0.667 & 0.422 & 0.199\\
1 & $10^{-3}$ & 1.530 (0.143) & 0.469 & 0.539 & 0.423 & 0.543\\
1 & $10^{-2}$ & 0.652 (0.031) & 0.466 & 0.483 & 0.424 & 0.845\\
1 & $10^{-1}$ & 0.486 (0.010) & 0.465 & 0.465 & 0.425 & 0.947\\
1 & $1$ & 0.443 (0.007) & 0.442 & 0.444 & 0.432 & 0.991\\
1.5 & $10^{-4}$ & 1.215 (0.071) & 0.582 & 0.586 & 0.509 & 0.670\\
1.5 & $10^{-3}$ & 1.110 (0.061) & 0.579 & 0.585 & 0.509 & 0.704\\
1.5 & $10^{-2}$ & 0.841 (0.035) & 0.575 & 0.580 & 0.509 & 0.808\\
1.5 & $10^{-1}$ & 0.611 (0.012) & 0.559 & 0.564 & 0.509 & 0.930\\
1.5 & $1$ & 0.536 (0.007) & 0.539 & 0.541 & 0.512 & 0.988\\\bottomrule
\end{tabular}\label{tab:a2}
\end{table}

第一次报告中的对应GCV-MA端点为：$p/n=1$ 时风险5.012至0.421、校准比0.227至1.017；$p/n=1.5$ 时风险1.328至0.545、校准比0.632至0.974。新种子结果并未复制逐个数值，但两个维数比的方向一致。由于候选集合随下界改变，风险下降不能单独归因于GCV估计改善；Oracle-MA风险在 $p/n=1$ 时从0.422至0.432，在 $p/n=1.5$ 时从0.509至0.512，显示缩小集合存在小幅近似代价。GCV-MA的收益远大于这部分Oracle代价。完整均值、标准误、原报告对照及配对差值见CSV。
\fig{A2_boundary_replication}{.93}{独立复验。上排为四种方法的平均超额风险及约95\% Monte Carlo均值区间，下排为各方法平均GCV分数除以平均总风险。横轴是候选下界的10为底对数；校准比1为均值层面的参照线。}

\section{唯一主要B研究问题及相对A的变化}
\textbf{主要问题：}能否仅利用训练数据内的五折折外预测误差选择GCV-MA的候选岭参数下界，使其在插值附近达到接近较优固定下界的预测风险，同时改善GCV风险校准？

相对A2，数据生成机制和GCV-MA选权器保持一致；改变的是候选集合的确定方式。A2事先固定下界并逐一诊断，B在每份训练数据上从五个下界中选一个，再用全部训练样本重新拟合。具体预实验对每个候选下界使用同一五折划分：每折用80个样本拟合GCV-MA，预测其余20个样本；平均100个折外平方误差后选最小者。固定 $10^{-4}$ 网格、固定下界1、单次80/20留出选界、同准则单一Ridge-GCV和LOO-MA构成基准。若真实参数和测试数据进入选界，将破坏问题的可实施性；程序将二者隔离于训练后评价阶段。

评价包括真实参数给出的超额风险均值及标准误、与固定网格的配对风险差、GCV分数与总风险均值比、不同下界的选择频数；同时记录Oracle-MA风险以衡量候选集合代价。B的预期假设是五折选界较完整网格降低风险，并接近固定下界1；不预设它在全部设计下优于固定1。

\section{B预实验：实际运行与假设判定}
预实验另用种子根20261023，仍取 $n=100$、Toeplitz协方差、信噪比1与两种维数比；每种维数比运行100份独立数据，共200份。五种下界及所有选界规则共享同一份数据和折分。此处没有使用A2复验种子来挑选规则；规则和基准的执行记录可从\texttt{Results/B\_pilot\_raw.csv}及\texttt{B\_pilot\_selections.csv}复算。
\begin{table}[!htbp]\centering\small
\caption{B预实验的平均超额风险（均值标准误）及校准比，每组100份数据。}
\begin{tabular}{clcc}\toprule
$p/n$ & 方法 & 超额风险 $\pm$ SE & 校准比\\\midrule
1 & 完整网格 $10^{-4}$ & $5.540\pm0.662$ & 0.202\\
1 & 单次留出选界 & $1.456\pm0.374$ & 0.574\\
1 & 五折选界 & $0.446\pm0.009$ & 0.984\\
1 & 固定下界1 & $0.441\pm0.009$ & 0.989\\
1.5 & 完整网格 $10^{-4}$ & $1.359\pm0.093$ & 0.618\\
1.5 & 单次留出选界 & $0.728\pm0.048$ & 0.856\\
1.5 & 五折选界 & $0.545\pm0.014$ & 0.968\\
1.5 & 固定下界1 & $0.526\pm0.009$ & 0.982\\\bottomrule
\end{tabular}\label{tab:b}
\end{table}

五折选界减完整网格的配对平均风险差，在 $p/n=1$ 为 $-5.094$（标准误0.662），在 $p/n=1.5$ 为 $-0.815$（标准误0.093）。减固定下界1的差值分别为0.005（标准误0.003）及0.019（标准误0.011），方向上仍略差；这两个配对差值的近似95\%区间都跨0。同数据上的完整网格Ridge-GCV风险为0.825和0.564，LOO-MA为0.650和0.575。五折规则在两组各100次中均有91次选择下界1，说明本预实验中接近固定1主要由频繁选到该下界驱动。单次留出会选到极小下界，因而风险尾部较重。\textbf{判定：}指定设置下的预实验支持相对完整网格改善风险和校准，尚不足以支持普遍优于或严格等价于固定1，更不能外推到其他信号强度和协方差。
\fig{B_pilot_risk}{.88}{B预实验的平均超额风险及约95\% Monte Carlo均值区间。每个维数比的四种规则使用同一100份数据；正式比较以配对差值CSV为准。}
\FloatBarrier

\section{B执行计划、工作量与其他探索方向}
拟将正式B设计扩展为 $n\in\{100,200\}$、$p/n\in\{0.95,1,1.05,1.5\}$、协方差\{独立、Toeplitz\}、信噪比\{1,5\} 的32个配置，每配置200次重复，共6,400份合成数据；各规则共享数据与折分。对五折选界、完整网格、固定下界1、单次留出、Ridge-GCV及LOO-MA报告配对风险与校准。固定下界1是当前设置的强基准；若五折选择在多数新配置中仍退化为固定1，应记录额外计算开销并缩小B的适用结论。本次200份预实验约运行18秒，据此粗略估计核心模拟约10分钟，图表、核验及写作需1--2个工作日。

其他可探索方向仅作为后续延伸：检查非高斯噪声与异方差下的选界稳定性；分析不同协方差谱和信号方向对候选集合Oracle风险的影响；研究混合平滑器的高维校准误差。这些方向不构成本次额外B主问题，也未在本报告中声称已验证。

\section{可复现性与限制}
\texttt{Code/a2\_core.py}由第一次A2实现复制并仅将随机种子根改为20261022；\texttt{run\_a2.py}筛取原配置中的10个boundary场景。B程序单列，分析脚本从原始CSV重新生成所有表和图。数值核验对首个A2重复重新运行，记录的12种方法超额风险最大绝对误差为0；18,000条A2方法记录的优化均成功，最大归一化一阶间隙为 $7.43\times10^{-6}$。上述核验不能代替对全部随机重复的形式证明。研究只覆盖指定合成数据和有限重复，没有真实数据验证、混合平滑器高维极限证明或普遍校准保证。代码、配置、原始记录、图表与本报告均随A2+B文件夹提交。

\begin{thebibliography}{99}
\bibitem[Dobriban and Wager(2018)]{dw}
Dobriban, Edgar; Wager, Stefan (2018). High-dimensional asymptotics of prediction: Ridge regression and classification. \textit{The Annals of Statistics}, 46, 247--279. \url{https://doi.org/10.1214/17-AOS1549}.
\bibitem[Yu et al.(2018)]{yu}
喻达磊、饶炜东、尹潇潇 (2018). 岭回归中基于广义交叉核实法的最优模型平均估计. \textit{系统科学与数学}, 38, 652--661. \url{https://doi.org/10.12341/jssms13408}.
\end{thebibliography}
\end{document}
